\paragraph{2.3. Разработка алгоритмов}
\subparagraph{2.3.1}
Лень, надо рисунок делать.

\subparagraph{2.3.2}
Перепишите процедуру MERGE так, чтобы в ней не использовались сигнальные значения.
Сигналом к остановке должен служить тот факт, что все элементы массива L или R скопированны обратно в массив A,
после чего в этот массив копируются элементы, оставшиеся в непустом массиве.\\
\\
\begin{tabular}{ll}
Merge($A, p, q, r$) & \\
\quad $n_1 = q - p + 1$ & \\
\quad $n_2 = r - q$ & \\ 
\quad Пусть $L[1 \ldots n_1]$ и $R[1 \ldots n_2]$ - новые массивы & \\
\quad for $i = 1$ to $n_1$ & \\
\quad \quad $L[i] = A[p + i - 1]$ & \\
\quad for $j = 1$ to $n_2$ & \\
\quad \quad $R[j] = A[q + j]$ & \\
\quad $i = 1$ & \\
\quad $j = 1$ & \\
\quad for $k = p$ to $r$ & \\
\quad \quad if $i > n_1$ & \\
\quad \quad \quad $A[k] = R[j]$ & \\
\quad \quad \quad $j = j + 1$ & \\
\quad \quad else if $j > n_2$ & \\
\quad \quad \quad $A[k] = L[i]$ & \\
\quad \quad \quad $i = i + 1$ & \\
\quad \quad else if $L[i] \leq R[j]$ & \\
\quad \quad \quad $A[k] = L[i]$ & \\
\quad \quad \quad $i = i + 1$ & \\
\quad \quad else & \\
\quad \quad \quad $A[k] = R[j]$ & \\
\quad \quad \quad $j = j + 1$ & \\
\end{tabular} \\

\subparagraph{2.3.3}
Воспользуйтесь методом математической индукции для доказательства того, что, когда $n$ является степенью 2, решением рекуррентного соотношения 
\[
T(n) = \begin{cases}
2 & \text{если } n = 2 \\
2T(n/2) + n & \text{если } n = 2^k, k > 1
\end{cases}
\] 
является $T(n) = n \lg n$.

\textbf{Доказательство:}

\textbf{База индукции:} $n = 2$.
\[
T(2) = 2, \quad 2 \lg 2 = 2 \cdot 1 = 2.
\]
Таким образом, $T(2) = 2 \lg 2$, база выполняется.

\textbf{Индукционный переход:}
Предположим, что для $n = 2^k$, где $k \ge 1$, выполняется $T(2^k) = 2^k \lg(2^k) = 2^k \cdot k$.
Докажем, что тогда формула верна для $n = 2^{k+1}$.

По рекуррентному соотношению:
\[
T(2^{k+1}) = 2T\left(\frac{2^{k+1}}{2}\right) + 2^{k+1} = 2T(2^k) + 2^{k+1}.
\]
Подставим индукционное предположение $T(2^k) = 2^k \cdot k$:
\[
T(2^{k+1}) = 2 \cdot (2^k \cdot k) + 2^{k+1} = 2^{k+1} \cdot k + 2^{k+1} = 2^{k+1}(k + 1).
\]
Заметим, что $\lg(2^{k+1}) = k + 1$, следовательно:
\[
T(2^{k+1}) = 2^{k+1} \lg(2^{k+1}) = n \lg n.
\]

Индукционный переход доказан.

\subparagraph{2.3.4}
Сортировку вставками можно представить в виде рекурсивной последовательности. 
Чтобы отсортировать массив $A[1 \ldots n]$, сначала нужно рекурсивно отсортировать массив $A[1 \ldots n-1]$, 
после чего в этот отсортированеный массив помещается элемент $A[n]$. 
Запишите рекурретное уравнение для времени работы этой рекурсивной версии сортировки вставкой.

\textbf{Рекурретное уравнение:}
\[
T(n) = \begin{cases}
1 & \text{если } n = 1 \\
T(n - 1) + n & \text{если } n > 1
\end{cases}
\] 

\subparagraph{2.3.5}
Возвращаясь к задаче поиска (см. упр. 2.1.3), нетрудно заметить, что если последовательность $A$ отсортирована,
то можно сравнить  значение среднего элемента этой последовательности с искомым значением $\upsilon$ и
сразу исключить половину последовательности из дальнейшего рассмотрения. \textbf{Бинарный поиск} (binary search) - это алгоритм,
в котором такая процедура повторяется неоднократно, что свякий раз приводит к уменьшению осташейся части последовательности в два раза.
Запишите псевдокод алгоритма бинраного поиска (либо итеративный, либо рекурсивный).
Докажите, что время работы этого алгоритма в наихудшем случае составялет $\Theta(\lg{n})$.

Полагаем массив $A \text{ содержит } k$.

\textbf{Псевдокод:}

\begin{tabular}{ll}
Binary-Search($A, k$) & \\
\quad $l = 0$ & \\
\quad $r = A.length + 1$ & \\
\quad while $l < r - 1$ & \\
\quad \quad $m = \lfloor (l + r)/2 \rfloor$ & \\
\quad \quad if $A[m] < k$ & \\ 
\quad \quad \quad $l = m$ & \\ 
\quad \quad else & \\
\quad \quad \quad $r = m$ & \\
\quad if $r \le A.length$ and $A[r] == k$ & \\
\quad \quad return $r$ & \\
\quad else & \\
\quad \quad return $NIL$ &
\end{tabular}

\textbf{Доказательство:}
В наихудшем случае искомый элемент отсутствует в массиве или находится на последней проверяемой позиции. 
На каждой итерации цикла область поиска уменьшается примерно вдвое:
после первой итерации остаётся не более $\lceil n/2 \rceil$ элементов, после второй — не более $\lceil n/4 \rceil$, и так далее.
Через $k$ итераций размер области поиска не превышает $\lceil n / 2^k \rceil$.
Процесс продолжается, пока область поиска не станет пустой, т.е. $n / 2^k < 1$.
Отсюда $k > \log_2 n$, то есть количество итераций равно $\lfloor \log_2 n \rfloor + 1 = \Theta(\log n)$.
Поскольку каждая итерация выполняется за $\Theta(1)$, общее время работы в наихудшем случае составляет $\Theta(\log n)$.

\subparagraph{2.3.6}
Заметим, что в цикле \textbf{while} в строках 5-7 процедуры Insertion-Sort в раделе 2.1 для сканирования (в обратном порядке) отсортированного помассива $A[1 \ldots j -1]$ используется линейный поиск. 
Можно ли использовать бинарный поиск (см. упр. 2.3.5) вместо линейного поиска, чтобы время работы этого алгоритма в наихудшем случае улучшилось и стало $\Theta(n\lg{n})$? 